\newlabel{fig:supp_av45_1}{{S1}{25}{\textbf {Comparison of real and synthetic AV45 uptake across the remaining brain regions (Part 1 of 2).} Values represent mean regional intensities from per-image min--max-normalized PET images (arbitrary units), rather than SUVR.\relax }{figure.caption.15}{}}
\newlabel{fig:supp_av45_2}{{S2}{26}{\textbf {Comparison of real and synthetic AV45 uptake across the remaining brain regions (Part 2 of 2).} Values represent mean regional intensities from per-image min--max-normalized PET images (arbitrary units), rather than SUVR.\relax }{figure.caption.16}{}}
\newlabel{fig:supp_fdg_1}{{S3}{27}{\textbf {Comparison of real and synthetic FDG uptake across the remaining brain regions (Part 1 of 2).} Values represent mean regional intensities from per-image min--max-normalized PET images (arbitrary units), rather than SUVR.\relax }{figure.caption.17}{}}
\newlabel{fig:supp_fdg_2}{{S4}{28}{\textbf {Comparison of real and synthetic FDG uptake across the remaining brain regions (Part 2 of 2).} Values represent mean regional intensities from per-image min--max-normalized PET images (arbitrary units), rather than SUVR.\relax }{figure.caption.18}{}}
\newlabel{fig:supp_disease_av45}{{S5}{29}{\textbf {Validation of disease-related AV45 uptake patterns across the remaining brain regions.} Violin plots display mean normalized regional uptake (arbitrary units) in Alzheimer’s disease (AD, red) and cognitively normal (CN, green) subjects for selected brain regions. Each region contains paired half-violins: the left half (solid fill) represents real PET data, and the right half (dashed fill) represents MRI-synthetic PET data, with consistent color coding between modalities. Solid lines above the violins indicate statistical significance for real PET, while dashed lines indicate significance for synthetic PET. Uptake values were derived from per-image min--max-normalized PET images without reference-region normalization.\relax }{figure.caption.19}{}}
\newlabel{fig:supp_disease_fdg}{{S6}{30}{\textbf {Validation of disease-related FDG uptake patterns across the remaining brain regions.} Violin plots display mean normalized regional uptake (arbitrary units) in Alzheimer’s disease (AD, red) and cognitively normal (CN, green) subjects for selected brain regions. Each region contains paired half-violins: the left half (solid fill) represents real PET data, and the right half (dashed fill) represents MRI-synthetic PET data, with consistent color coding between modalities. Solid lines above the violins indicate statistical significance for real PET, while dashed lines indicate significance for synthetic PET. Uptake values were derived from per-image min--max-normalized PET images without reference-region normalization.\relax }{figure.caption.20}{}}
\newlabel{tab:supp_adni}{{S1}{31}{\textbf {Demographic characteristics and prepared clinical inputs of ADNI.} Data are mean $\pm $ sample s.d., except counts. Demographic $P$ values use one-way ANOVA (age, weight) or a chi-square test (sex). Clinical summaries combine observed and diagnostic-group-conditioned filled scores; their group-comparison $P$ values are omitted because the distributions are not independent clinical evidence. Source missingness is reported in Table~\ref {tab:clinical_missingness}. CDR: Clinical Dementia Rating; MMSE: Mini-Mental State Examination.\relax }{table.caption.21}{}}
\newlabel{tab:supp_oasis}{{S2}{31}{\textbf {Demographic characteristics of OASIS-3.}\relax }{table.caption.22}{}}
\newlabel{tab:supp_software}{{S3}{32}{\textbf {Historical core software dependencies and reported versions for DIReCT$++$.} This retained environment list describes the historical configuration; reconstructed training and the timing run used separately recorded environments. In particular, the RTX 4090 timing run used PyTorch 2.5.1+cu121, as stated in Methods. The table is not a single environment specification for every experiment.\relax }{table.caption.23}{}}
\newlabel{tab:clinical_missingness}{{S4}{32}{\textbf {Missing clinical scores in the source records before historical filling.} Each cell gives missing count / group size (percentage); missingness is defined relative to the first source CSV record selected for each participant, not all ADNI visits.\relax }{table.caption.24}{}}
\newlabel{tab:reproducibility}{{S5}{33}{\textbf {Reproducibility settings and scope of the revised analyses.}\relax }{table.caption.25}{}}
\newlabel{tab:supp_tau_adapter}{{S6}{34}{\textbf {Adapter-level sensitivity analysis of the cross-tracer soft target.} Descriptive means across two folds, using 1,856 QC-filtered participants. Only training-fold PET features were used for fitting; held-out features were evaluated after optimization. No cross-term means removal of the penalty ($\lambda _{\mathrm {cross}}=0$), not a target of zero. No generator was retrained or evaluated in this analysis.\relax }{table.caption.27}{}}
\newlabel{tab:regional_agreement}{{S7}{35}{\textbf {Regional normalized uptake agreement and descriptive regression for the 16 representative tracer--ROI combinations in Figure~\ref {fig:4}A,B.} Overall cohorts from the author-designated historical ROI CSVs; N can vary by ROI.\relax }{table.caption.28}{}}
\newlabel{tab:regional_effect_sizes}{{S8}{36}{\textbf {AD--CN regional effect sizes and BH-FDR-adjusted comparisons for Figure~\ref {fig:4}C--F.} Cohen's $d$ is AD minus CN divided by the pooled within-group sample s.d.\relax }{table.caption.29}{}}
\newlabel{tab:whole_volume_generation}{{S9}{37}{Whole-volume OOF generation metrics (mean $\pm $ sample s.d.)\relax }{table.S9}{}}
\newlabel{tab:paired_generation}{{S10}{38}{Paired whole-volume generation comparisons\relax }{table.S10}{}}
\newlabel{tab:r3_11_masked_summary}{{S11}{40}{Mask-restricted quality of frozen two-fold OOF PET predictions.\relax }{table.caption.30}{}}
\newlabel{tab:r3_11_masked_paired}{{S12}{40}{Paired mask-restricted PET image-quality comparisons.\relax }{table.caption.31}{}}
\newlabel{fig:cases_fdg}{{S7}{41}{\textbf {FDG synthesis in median, 10th-percentile and worst QC-passed cases.} Rows follow the frozen whole-volume SSIM selection; columns show MRI, real PET, synthetic PET and absolute-error maps. PET intensities and errors use the displayed [0,1] normalized-intensity scales. Selection is within the existing QC-passed set, without new exclusions or visual replacement; rankings use whole-volume SSIM rather than masked SSIM.\relax }{figure.caption.32}{}}
\newlabel{fig:cases_av45}{{S8}{42}{\textbf {AV45 synthesis in median, 10th-percentile and worst QC-passed cases.} Rows follow the frozen whole-volume SSIM selection; columns show MRI, real PET, synthetic PET and absolute-error maps. PET intensities and errors use the displayed [0,1] normalized-intensity scales. Selection is within the existing QC-passed set, without new exclusions or visual replacement; rankings use whole-volume SSIM rather than masked SSIM.\relax }{figure.caption.33}{}}
\newlabel{tab:classification_summary}{{S13}{43}{Three-fold classification results: mean $\pm $ sample s.d.; [conditional 95\% CI].\relax }{table.S13}{}}
\newlabel{tab:classification_comparisons}{{S14}{45}{Paired imaging classification differences\relax }{table.S14}{}}
\newlabel{tab:classification_clinical_comparisons}{{S15}{46}{Paired comparisons with clinical baselines\relax }{table.S15}{}}
\newlabel{tab:classification_delong}{{S16}{48}{Paired DeLong comparisons by outer fold\relax }{table.S16}{}}
\newlabel{fig:classification_roc}{{S9}{50}{\textbf {Cross-sectional classification ROC curves for AD/CN, EMCI+LMCI/CN and EMCI/LMCI (left to right).} Curves average three fold-specific TPRs at a fixed FPR grid; shaded regions are pointwise conditional 95\% bootstrap bands, not simultaneous confidence bands. Legend AUC values are equally weighted empirical fold AUC means. Comparisons include MRI, MRI+Real both, MRI+Syn both, Clinical-only and MRI+Clinical. The clinical fields include diagnosis-related scores and historical group-conditioned fills; high clinical AUC is not evidence of independent synthetic-PET benefit.\relax }{figure.caption.34}{}}
\newlabel{tab:e1_score_removal}{{S17}{50}{\textbf {E1 fixed-weight score-removal sensitivity.} Differences are No-Scores minus Full, except the nonnegative output difference between predictions. This is a distinct E1 Full checkpoint family, not the reconstructed Full used in the generator-comparison tables.\relax }{table.caption.35}{}}
\newlabel{tab:r2-16-ad-cn}{{S18}{51}{Age- and sex-stratified classification performance: AD vs CN.\relax }{table.caption.36}{}}
\newlabel{tab:r2-16-mci-cn}{{S19}{51}{Age- and sex-stratified classification performance: EMCI+LMCI vs CN.\relax }{table.caption.37}{}}
\newlabel{tab:r2-16-emci-lmci}{{S20}{51}{Age- and sex-stratified classification performance: EMCI vs LMCI.\relax }{table.caption.38}{}}
